%document class
\documentclass[10pt,oneside]{book}
%%%% Page Info + Commands %%%%%{

%packages
\usepackage{pgfplots}
\pgfplotsset{compat=1.18}
\usepackage{amsmath}
\usepackage{amsfonts}
\usepackage{amssymb}
\usepackage{amsthm,letltxmacro,changepage}
\usepackage{enumerate}
\usepackage{enumitem}
\usepackage{changepage}

%This will shrink the page
\newcommand{\smallpage}{  \setlength \oddsidemargin{.5in}
            \setlength \textwidth{5in}
            \setlength \topmargin{0in}
            \setlength \textheight{9in}}


% Commands for sets of numbers
\newcommand{\Z}{\mathbb{Z}}\newcommand{\R}{\mathbb{R}}\newcommand{\N}{\mathbb{N}}

\newcommand{\n}{\vspace{0.15cm}\\}
\newcommand{\en}{\vspace{0.25cm}\\}

\newcommand{\gimage}[3]{
\begin{figure}[h]   
    \centering          
    \includegraphics[width=#2\textwidth]{#1}  
    \caption{#3}
    \label{fig:#1}
\end{figure}\\
}

\newcommand{\centerit}[1]{{\begin{center} #1 \end{center}}}
\newcommand{\ul}[1]{\underline{#1}}

%%%%%%%% command for graphics %%%%%%%%%%%%%
\usepackage{fancyhdr}
%}

%%%% Page Setup %%%%%%%
\smallpage\sloppy
\pagestyle{fancy}
\lhead{\large{\textbf{Problem Set 2 - Gordie's Novak}}}
%\rfoot{\thepage/\pageref{LastPage} }
\setlength{\headheight}{10pt}


% Proof writing
\LetLtxMacro\oldproof\proof
\let\endoldproof\endproof
\renewenvironment{proof}[1][\proofname]
  {\vspace{0.2cm}\begin{adjustwidth}{2em}{2em}
   \oldproof[#1]}
  {\endoldproof
   \end{adjustwidth}}
\begin{document}

\subsection*{Section 2.3}
\begin{enumerate}
    %
    % Question 2
    %
    \item[2.] Given integers $a, b, e, d$ verify the following:
    \begin{enumerate}
        \item[(b)] If $a \mid b$ and $a \mid c$, then $a^2 \mid bc$.\en
        We will show that because $a$ divides both $b$ and $c$, that $a^2$ will divide $b\cdot c$. We will do this by using the definition of ``divides'' to our advantage.
        \begin{proof}
            We want to show that if $a \mid b$ and $a \mid c$, then $a^2 \mid bc$. \n
            Consider that because $a\mid b$, then $\exists n\in \Z$ such that $n\cdot a = b$. The same applies to $a\mid c$, such that $\exists m\in\Z$ such that $m\cdot a = c$.\en
            Thus, we can represent $bc$ as:
            \begin{align*}
                bc &= an\cdot am\\
                &= nm\cdot a^2
            \end{align*}
            Let $\ell \equiv nm$, $\ell \in \Z$. Hence, via the definition of ``divides'', because $\ell a^2 = bc$, $a^2\mid bc$, as desired. 
        \end{proof}
    \end{enumerate}
    %
    % Question 3
    %
    \item [3.]Prove or disprove: If $a \mid (b + c)$, then either $a \mid b $ or $a \mid c$.\en
    I will disprove this statement with a simple counterexample.
    \begin{proof}
        We want to show by providing a counterexample that if $a\mid (b+c)$, that does not imply that either $a\mid b$ or $a\mid c$. \en
        Let $a = 2$, $b = 1$, and $c = 1$. We have that $2\mid (1+1) = 2$. \en
        However, note that because $a > b,c$, that $a\nmid b, c$. Hence, if $a\mid (b+c)$, that does not imply that either $a\mid b$ or $a\mid c$, as desired.  
    \end{proof}
    %
    % Question 12
    %
    \item [12.] Prove that, for a positive integer $n$ and any integer $a,\; \text{gcd}(a, a+ n)$ divides $n$.\n
    We will prove this by selecting an arbitrary $m$ that helps to satisfy the condition of this problem.
    \begin{proof}
        Let $a, n\in \Z$. We want to show that for $m = \text{gcd}(a,\;a+n)$, $m$ divides $n$. \n
        Consider that because $m$ is a divisor of $a$ and $a+n$, $\exists x, y\in \Z$ such that $x\cdot m = a$ and $y\cdot m = a+n$.\n
        Now, we write that:
        \begin{align*}
            ym &= a+n \\
            ym&= xm + n\\
            (y-x)m &= n
        \end{align*}
        Let $\ell \equiv (y-x)\in \Z$. Hence, because $\exists \ell \in \Z$ such that $\ell m = n$, $m\mid n$. Because $m = \text{gcd}(a,a+n)$, we equivilantly state that the greatest common divisor of $a$ and $a+n$ divides $n$, as desired. 
    \end{proof}
    \pagebreak
    %
    % Question 12
    %
    \item [14.] For any integer $a$, show the following: $\text{gcd}(2a + 1 , 9a + 4) = 1$.\n
    I will show in this proof that there exists a linear combination of $2a +1$ and $9a +4$ such that said combination equals one.
    \begin{proof}
        Let $a$ be an integer. We want to show that: 
        \[\text{gcd}(2a + 1 , 9a + 4) = 1.\]
        Consider the following linear combination of $2a+1$ and $9a+4$:
        \[(9\cdot(2a+1)) + (-2\cdot(9a+4)) = 1\]
        Hence, via \textbf{Theorem 2.4}, because there exists a linear combination of both terms that equals 1, $2a+1$ and $9a+4$ are relatively prime, and thus their greatest common denominator is 1, as desired.
    \end{proof}
    \item [20b. ]Confirm the following properties of the greatest common divisor: If $\text{gcd}(a, b) = 1$, and $c \mid a$, then $\text{gcd}(b, c) = 1$.\n
    In this proof I will show with the help \textbf{Theorem 2.4} that $b$ and $c$ are relatively prime, and thus their gcd is 1. 
    \begin{proof}
        Let $a,b$ and $c$ be integers such that $\text{gcd}(a,b) = 1$ and $c\mid a$. We want to show that $\text{gcd}(b,c) = 1$.\n
        Via \textbf{Theorem 2.4}, because $a$ and $b$ are relatively prime, let $x,y\in \Z$ such that:
        \[xa + yb = 1\]
        Furthermore, via the properties of divisible numbers, because $c\mid a$, let $\ell\in\Z$ such that $\ell c = a$. Substituting $\ell c$ into the above equation gives:
        \[x\ell c + yb = 1\]
        Let $z \equiv x\ell\in \Z$. Thus, because $zc + yb = 1$, $c$ and  $b$ are relatively prime, and hence their greatest common denominator is one, as desired. 
    \end{proof}
\end{enumerate}
\subsection*{Section 2.4}
\begin{enumerate}
    \item Find gcd(143, 227) and gcd(306, 657).\n
    I will use a computer program implemented with the Euclidean algorithm to solve this problem that I have created. The primary snippet I am using is displayed below. For the complete implementation, refer to the \texttt{gcd.cpp} file in my submission page.
\begin{verbatim}
int a, b;
while (b != 0) {
    const int r = a % b;
    if (r == 0)
        break;
    a = b;
    b = r;
}
\end{verbatim}
    Running the program yields the following results:
    \begin{enumerate}
        \item gcd$(143, 227) = 1$
        \item gcd$(306,657) = 9$
    \end{enumerate}
    \pagebreak
    However, much to my personal grumblings, and upon reflection, I realize that I'll just do it with the Euclidean algorithm to make sure I get full credit on this assignment. (insert sad face emoji here)
    \begin{enumerate}
        \item gcd(143, 227)
        \begin{align*}
            227 &= 1(143) + 84\\
            143 &= 1(84) + 59\\
            84 &= 1(59) + 25\\
            59 &= 2(25) + 9 \\
            25 &= 2(9) + 7 \\
            9 &= 1(7) + 2 \\
            7 &= 3(2) + 1 \\
            2 &= 2(1) + 0
        \end{align*}
        We look at the second to last remainder and see that our gcd is 1.
        \item  gcd(306, 657)
        \begin{align*}
            657 &=  2(306) + 45\\
            306 &= 6(45) + 36 \\
            45 &= 1(36) + 9 \\
            36 &= 4(9) + 0
        \end{align*}
        We look at the second to last remainder and see that our gcd is 9.
    \end{enumerate}

\end{enumerate}

\end{document}
